\documentclass[../../../include/open-logic-section]{subfiles}

\begin{document}

\olfileid{sth}{spine}{recursion}
\olsection{Teorema-teorema Rekursi Transfinit}

Nanging luwih dhisik kita kudu mesthekake yen
\olref[valpha]{defValphas} pancen definisi kang sah.
Luwih umum, kita kudu mbuktekake yen netepake
sawijining obyek liwat rekursi transfinit bisa
ditindakake. Iki tujuan bagean iki.

\emph{Pangeling: materi iki rada angel}. Nanging
gagasan utamane prasaja: Induksi Transfinit bebarengan
karo Panggantian njamin sahé sawetara wujud rekursi
transfinit.\footnote{Pangeling: kabeh rumus lan term
bisa nduweni parameter (kajaba yen kanthi cetha
kasebut kosok baline).}

\begin{defn}
	Upamane $\tau(x)$ term, $f$ fungsi, lan $\alpha$
	ordinal. Kita ngarani $f$ \emph{aproksimasi-$\alpha$}
	kanggo $\tau$ yen lan mung yen $\dom{f} = \alpha$
	lan $(\forall \beta \in \alpha)f(\beta)
	= \tau(\funrestrictionto{f}{\beta})$.
\end{defn}

\begin{lem}[Rekursi Winates]\ollabel{transrecursionfun}
Kanggo sembarang term $\tau(x)$ lan sembarang
ordinal $\alpha$, ana aproksimasi-$\alpha$ kang
tunggal kanggo $\tau$.
\end{lem}
\begin{proof}
Kita bakal nuduhake yen kanggo saben
$\gamma \leq \alpha$ ana aproksimasi-$\gamma$
kang tunggal.

Pisanan kita buktekake tunggale. Upamane $g$ lan $h$
iku aproksimasi-$\gamma$ lan aproksimasi-$\delta$.
Induksi transfinit ing argumen-argumene nuduhake
$g(\beta) = h(\beta)$ kanggo saben
$\beta \in \dom{g} \cap \dom{h}
= \gamma \cap \delta = \min(\gamma, \delta)$.
Mula aproksimasi kasebut tunggal (yen ana), lan
padha ing kabeh nilai kang dienggo bebarengan.

Kanggo mbuktekake anane, saiki kita nggunakake
induksi transfinit prasaja
(\olref[ordinals][opps]{simpletransrecursion})
ing ordinal $\delta \leq \alpha$.

Fungsi kosong cetha dadi aproksimasi-$\emptyset$.

Yen $g$ aproksimasi-$\gamma$, mula
$g \cup \{\tuple{\gamma, \tau(g)}\}$ dadi
aproksimasi-$\ordsucc{\gamma}$.

Yen $\gamma$ ordinal limit lan $g_\delta$
aproksimasi-$\delta$ kanggo saben $\delta < \gamma$,
amarga saben aproksimasi tunggal, Panggantian bisa
nglumpukake kabeh aproksimasi iki dadi himpunan.
Upamane $g = \bigcup_{\delta \in \gamma} g_\delta$.
Iki fungsi, amarga kabeh $g_\delta$ sarujuk
ing nilai-nilai kang tumpang tindih. Yen
$\delta \in \gamma$, mula
$g(\delta) = g_{\ordsucc{\delta}}(\delta)
= \tau(\funrestrictionto{g_{\ordsucc{\delta}}}{\delta})
= \tau(\funrestrictionto{g}{\delta})$.

Iki ngrampungake bukti kanthi induksi transfinit.
\end{proof}

Yen kita netepake \emph{term} tinimbang fungsi,
wates $\alpha$ ing asil sadurunge bisa dicopot.
Ing pratelan lan bukti asil sabanjure, yen $\sigma$
term, kita netepake
$\funrestrictionto{\sigma}{\alpha}
= \Setabs{\tuple{\beta, \sigma(\beta)}}{\beta \in \alpha}$.

\begin{thm}[Rekursi Umum]\ollabel{transrecursionschema}
Kanggo sembarang term $\tau(x)$, kita bisa netepake
kanthi gamblang sawijining term $\sigma(x)$ kang
nyukupi $\sigma(\alpha)
= \tau(\funrestrictionto{\sigma}{\alpha})$
kanggo sembarang ordinal~$\alpha$.
\end{thm}

\begin{proof}
Kanggo saben $\alpha$, miturut
\olref{transrecursionfun} ana aproksimasi-$\alpha$
kang tunggal, yaiku $f_\alpha$, kanggo~$\tau$.
Temtokna $\sigma(\alpha)$ minangka
$f_{\ordsucc{\alpha}}(\alpha)$. Saiki:
	\begin{align*}
		\sigma(\alpha) &= 
		f_{\ordsucc{\alpha}}(\alpha) \\&= 
		\tau(\funrestrictionto{f_{\ordsucc{\alpha}}}{\alpha}) \\&= 
		\tau(\Setabs{\tuple{\beta, f_{\ordsucc{\alpha}}(\beta)}}{\beta \in \alpha}) \\&= 
		\tau(\Setabs{\tuple{\beta, f_{\ordsucc{\beta}}(\beta)}}{\beta \in \alpha})\\&=
		\tau(\funrestrictionto{\sigma}{\alpha})
	\end{align*}
ing kene $f_{\ordsucc{\beta}}(\beta)
= f_{\ordsucc{\alpha}}(\beta)$ kanggo saben
$\beta < \alpha$, kaya ing
\olref{transrecursionfun}.
\end{proof}
\noindent 
Elinga yen \olref{transrecursionschema} iku
\emph{skema}. Kang wigati, kita ora bisa nganggep
$\sigma$ minangka fungsi kang awujud \emph{himpunan},
amarga yen mangkono $\dom{\sigma}$ bakal dadi
himpunan kabeh ordinal, nentang Paradoks Burali-Forti
(\olref[ordinals][basic]{buraliforti}).

Nanging isih kudu dituduhake yen
\olref{transrecursionschema} mbenerake definisi
himpunan-himpunan $V_\alpha$. Iki mbokmenawa ora
langsung katon; nanging bakal cetha liwat wujud
rekursi transfinit pungkasan kang prasaja.

\begin{thm}[Rekursi Prasaja]\ollabel{simplerecursionschema}
Kanggo sembarang term $\tau(x)$ lan $\theta(x)$
uga sembarang himpunan $A$, kita bisa netepake
kanthi gamblang term $\sigma(x)$ kang nyukupi:
\begin{align*}
	\sigma(\emptyset) &= A\\
	\sigma(\ordsucc{\alpha}) &= \tau(\sigma(\alpha)) &&
		\text{kanggo sembarang ordinal }\alpha\\
	\sigma(\alpha) &= \theta(\ran{\funrestrictionto{\sigma}{\alpha}})&&
	\text{yen }\alpha\text{ ordinal limit}
\end{align*}
\end{thm}

\begin{proof}
Pisanan kita netepake term $\xi(x)$ mangkene:
%t $\xi(x) = A$ yen $x$ dudu fungsi kanthi domain ordinal utawa fungsi kosong;
%yen ora, temtokna:
\[
	\xi(x) = 
	\begin{cases}
			A &\text{yen $x$ dudu fungsi kang}\\
			&\hspace{1em}\text{domaine ordinal utawa fungsi iku kosong; yen ora:}\\
			\tau(x(\alpha)) & \text{yen $\dom{x} = \ordsucc{\alpha}$}\\
			\theta(\ran{x}) & \text{yen $\dom{x}$ ordinal limit}
	\end{cases}
\]
Miturut \olref{transrecursionschema}, ana term
$\sigma(x)$ kang nyukupi
$\sigma(\alpha) = \xi(\funrestrictionto{\sigma}{\alpha})$
kanggo saben ordinal $\alpha$; kajaba iku,
$\funrestrictionto{\sigma}{\alpha}$ iku fungsi
kanthi domain $\alpha$. Kita nuduhake yen $\sigma$
nduweni sipat-sipat kang dibutuhake nganggo induksi
transfinit prasaja
(\olref[ordinals][opps]{simpletransrecursion}).

Pisanan, $\sigma(\emptyset) = \xi(\emptyset) = A$.

Sabanjure,
$\sigma(\ordsucc{\alpha}) = \xi(\funrestrictionto{\sigma}{\ordsucc{\alpha}}) = \tau(\funrestrictionto{\sigma}{\ordsucc{\alpha}}(\alpha)) = \tau(\sigma(\alpha))$.

Pungkasan,
$\sigma(\alpha) = \xi(\funrestrictionto{\sigma}{\alpha}) = \theta(\ran{\funrestrictionto{\sigma}{\alpha}})$,
yen $\alpha$ ordinal limit.
\end{proof}
\noindent
Saiki, kanggo mbenerake \olref[valpha]{defValphas},
pilihen $A = \emptyset$, $\tau(x) = \Pow{x}$,
lan $\theta(x) = \bigcup x$. Pungkasane definisi
himpunan-himpunan $V_\alpha$ iku sah!{}


\end{document}
